%% kapa: exam=1 round=35 year=2024 period=1 form=A subject=civil_law subject_ko=민법 seq=14 number=14
%% source: Q-Net 감정평가사 기출문제 (cst00309) artlSeq=5213229 file=raw/1008/5213229/1교시 문제지 원본.pdf
%% conversion: kordoc 4.24.51 -> tools/gen1.py (auto); review=proofread-sonnet; flags=none
\begin{question}{14}
乙은 대리권 없이 甲을 위하여 甲소유의 X토지를 丙에게 매도하였다. 이에 관한 설명으로 옳지 않은 것은? (다툼이 있으면 판례에 따름)
\begin{choices}
  \item 乙이 丙으로부터 받은 매매대금을 甲이 수령한 경우, 특별한 사정이 없는 한 甲은 위 매매계약을 추인한 것으로 본다.
  \item 甲이 乙을 상대로 위 매매계약의 추인을 한 경우, 그 사실을 丙이 안 때에는 甲은 丙에게 추인의 효력을 주장할 수 있다.
  \item 甲을 단독상속한 乙이 자신의 매매행위가 무효임을 주장하는 것은 신의칙에 반하여 허용되지 않는다.
  \item 丙이 甲에게 기간을 정하여 그 추인 여부의 확답을 최고하였으나 甲이 기간 내에 확답을 발송하지 않으면 추인은 거절한 것으로 본다.
  \item 甲이 추인을 하더라도 丙은 乙을 상대로 무권대리인의 책임에 따른 손해배상을 청구할 수 있다.
\end{choices}
\end{question}
\answer{5}
